%! The Language of Variation — Unit 1, Lesson 1.1 — Experience & Formalize
%! (Activity · QuickNotes · Application) (Answer Key)
%! Directory stays `experience/` (a build identifier in shared/lesson.mk); the label
%! students and teachers read is "Experience & Formalize".
% EFFL component, Math Medic sizing: 12pt body, OPEN white space after every question.
% \answerspace{H}{} reserves exactly H in the blank; the key passes the answer as the second
% argument, so the two files paginate identically by construction. Keep the macro
% byte-identical in the blank and the key. \boxguard counts here are 12pt-relative (~14-16).
% Page budget: Activity <=2pp · QuickNotes 1/2pp · Application 1/2-1pp.
% Check Your Understanding is NOT in this component: those items were moved to the end of
% the packet and are now the lesson's SCORED homework (see ../homework/).
%
% SPOILER RULE: the Activity never says categorical, quantitative, numerical, or measurement
% variable. Students invent their own two piles and name them in their own words ("amounts",
% "which kind"); the debrief attaches the formal terms in QuickNotes. `Individual', `variable',
% and `value' WERE formalized in Lesson 1.0, so they are prior knowledge here, not spoilers.
%
% GENERATED FROM experience/main.tex — this file is the blank with answers filled in, so the
% \answerspace heights, the \begin{work} block, and every \boxguard are identical to it.
% NO teachernote here — teacher prose lives in the lesson plan.
\documentclass[12pt]{article}
\usepackage{apstats-article}
\usepackage{apstats-key}   % pulls in -boxes; do NOT also load -boxes
\newcommand{\answerspace}[2]{\par\nopagebreak\noindent\begin{minipage}[t][#1][t]{\linewidth}%
  \color{keyred}\bfseries #2\end{minipage}\par}

\begin{document}

\pageheader{Unit 1, Lesson 1.1}{Experience \& Formalize: The Trailhead Register --- Answer Key}
% NAMESTRIP: no \namedateperiod here — mirrors the blank.

\vspace{0.02in}

% ── Part 1: the Activity ─────────────────────────────────────────────────────
% Scenario 1 makes the groups INVENT the two piles and name them themselves — no bins are
% handed to them. Scenario 2 breaks the naive rule they will have written ("has digits in it")
% with three digit-written labels and one genuine count, and carries the CRUX in 2b: the
% average of two permit numbers computes fine and means nothing.
\begin{headlinebox}{sky}
\textbf{The Trailhead Register.} Every party that hikes at Ross Hollow State Park signs the
register at the trailhead and signs out again on the way home. Work through the two problems
below with your group using only what you already know, and be ready to explain \emph{how you
know} each answer \textbf{from the register itself}.
\end{headlinebox}

\vspace{0.04in}

\begin{scenariobox}[1. Two Piles]{navy}
\small
\renewcommand{\arraystretch}{1.12}
\begin{tabularx}{\linewidth}{@{} l Y Y Y Y Y @{}}
\rowcolor{sky}
\textbf{Party} & \textbf{Trail} & \textbf{Group size} & \textbf{Time out (min)} &
\textbf{Dog along} & \textbf{Weather} \\
\hline
Alvarez  & Ridge Loop   & 4 &  95 & Yes & Sunny  \\
Boone    & Cedar Falls  & 2 & 140 & No  & Cloudy \\
Chen     & Ridge Loop   & 6 &  80 & Yes & Sunny  \\
Delgado  & Summit Trail & 3 & 215 & No  & Rain   \\
Egan     & Cedar Falls  & 2 & 155 & Yes & Cloudy \\
\hline
\end{tabularx}

\normalsize
\begin{enumerate}[label=\textbf{\alph*.}, leftmargin=1.8em, itemsep=6pt]
  \item One \textbf{row} of this register describes \ans{one hiking party}.
  \item Sort the five recorded headings into \textbf{two piles of your own}. Write each
        heading under the pile where you would put it.

        \begin{minipage}[t]{0.47\linewidth}
        \textbf{Pile 1}
        \answerspace{1.4cm}{Group size\\ Time out (min)\\[2pt]
        \emph{(amounts)}}
        \end{minipage}\hfill
        \begin{minipage}[t]{0.47\linewidth}
        \textbf{Pile 2}
        \answerspace{1.4cm}{Trail \quad Dog along \quad Weather\\[2pt]
        \emph{(which kind / which group)}}
        \end{minipage}
  \item Now name your piles. Finish this sentence \emph{twice}, once for each pile:
        \emph{``A heading belongs in this pile when \ldots{}''}
        \answerspace{2.0cm}{Pile 1: \ldots{} the entry is an amount of something you could
        measure or count, so the entries can be compared as sizes.\\[2pt]
        Pile 2: \ldots{} the entry names which group the party falls into, not how much of
        anything it has.}
  \item For \textbf{one} of your piles, adding up all five entries in a column and dividing by 5
        gives the ranger something useful. For the other it gives nonsense. Say which is which,
        write down the nonsense you would get on \textbf{Weather}, and say what you would do
        with that column instead to summarize it.
        \answerspace{2.4cm}{Pile 1 averages usefully: the mean time out is
        $(95+140+80+215+155)/5 = 137$ minutes. Pile 2 does not --- Sunny $+$ Cloudy $+$ Rain
        cannot be added at all, so there is no total to divide. Summarize Pile 2 by counting
        instead: 3 of the 5 parties (60\%) hiked with a dog.}
\end{enumerate}
\end{scenariobox}

\boxguard[16]
\begin{scenariobox}[2. Four Columns Written in Digits]{navy}
Next season the ranger adds four more columns. Here are the first three rows.

\vspace{4pt}
\small
\renewcommand{\arraystretch}{1.12}
\begin{tabularx}{\linewidth}{@{} l Y Y Y Y @{}}
\rowcolor{sky}
\textbf{Party} & \textbf{Permit \#} & \textbf{Parking spot} & \textbf{Home ZIP} &
\textbf{Water bottles} \\
\hline
Alvarez & 2041 & 17 & 80211 & 6 \\
Boone   & 6318 & 04 & 80304 & 2 \\
Delgado & 1150 & 11 & 80211 & 3 \\
\hline
\end{tabularx}

\normalsize
\begin{enumerate}[label=\textbf{\alph*.}, leftmargin=1.8em, itemsep=6pt]
  \item Every one of these four columns is written entirely in digits. Put each heading into
        \textbf{one of your two piles} from problem 1.

        \vspace{2pt}

        \textbf{Permit \#} \ans{Pile 2} \hfill \textbf{Parking spot} \ans{Pile 2}

        \vspace{4pt}
        \textbf{Home ZIP} \ans{Pile 2} \hfill \textbf{Water bottles} \ans{Pile 1}
  \item The Alvarez and Boone parties had permit numbers 2041 and 6318, and carried 6 and 2
        water bottles. Compute both averages.
        \begin{work}
        \text{average permit \#} &= \frac{2041 + 6318}{2} = 4179.5 \\[4pt]
        \text{average water bottles} &= \frac{6 + 2}{2} = 4
        \end{work}
        One of those two answers tells the ranger something real about the two parties and one
        does not. Which is which, and what exactly goes wrong with the other?
        \answerspace{1.8cm}{The bottle average is real --- between them the two parties carried
        4 bottles each. The permit average is nonsense: 4179.5 is not ``how much permit''
        anybody had. A permit number labels a party, so arithmetic on it answers nothing.}
  \item So being written in digits is not what decides the pile. Finish the rule in your own
        words: \emph{``A heading belongs in the pile you can average only when \ldots{}''}
        \answerspace{1.6cm}{\ldots{} the entries record \textbf{how much} or \textbf{how many}
        of something. If they only say \emph{which one} an individual is, the digits are a name.}
  \item The season report says \textbf{``62\% of parties hiked with a dog.''} That number came
        out of the \emph{Dog along} column. Does it move that column into the average pile?
        \answerspace{2.0cm}{No. The 62\% is a summary the ranger \emph{built} by counting Yes
        entries; it is not what is written in the column. Every entry there is still Yes or No
        --- a group label --- so the column stays in Pile 2. Counting a group column always
        produces numbers, and that never changes what the column records.}
\end{enumerate}
\end{scenariobox}

% ── Part 2: QuickNotes (the debrief fills this) — target: half a page ────────
\boxguard[14]
\begin{tcolorbox}[enhanced, colback=sky, colframe=navy, boxrule=0.8pt, arc=2mm,
    left=3mm, right=3mm, top=1.5mm, bottom=1.5mm, breakable,
    fonttitle=\bfseries\small\color{white}, title={QuickNotes: Two Kinds of Variable},
    attach boxed title to top left={yshift=-2mm, xshift=4mm},
    boxed title style={colback=navy, colframe=navy, arc=1.5mm}]
\small
\begin{minipage}[t]{0.34\linewidth}
\vspace{2pt}
\renewcommand{\arraystretch}{1.15}
\begin{tabular}{@{}l l@{}}
\multicolumn{2}{@{}l}{\footnotesize\textbf{\color{navy}Trail} \quad \textbf{\color{navy}Time (min)}}\\
\hline
Ridge Loop   & 95  \\
Cedar Falls  & 140 \\
Summit Trail & 215 \\
\hline
\end{tabular}\\[3pt]
{\footnotesize Left column: \textbf{category}\\
\textbf{names}. Right column:\\
a \textbf{measured amount}.\\
Same five parties, two\\
different kinds of column.}
\end{minipage}\hfill
\begin{minipage}[t]{0.62\linewidth}
\vspace{-4pt}
\begin{itemize}[leftmargin=1.1em, itemsep=5pt]
  \item A \textbf{variable} is a characteristic that \ans{changes} from one individual to
        another.
  \item A \textbf{categorical variable} takes values that are \ans{category} names or
        group \ans{labels}.
  \item A \textbf{quantitative variable} takes \ans{numerical} values for a measured or
        \ans{counted} quantity.
  \item \textbf{The test is not ``does it have digits.''} Ask: is the value an
        \ans{amount} of something? If adding or averaging the values is
        \ans{meaningful}, it is quantitative. A ZIP code is all digits and is
        \ans{categorical}.
  \item Counting a categorical variable \emph{produces} numbers --- counts and percents.
        The \ans{summary} is numeric; the \ans{variable} is still categorical.
\end{itemize}
\end{minipage}
\end{tcolorbox}

% ── Part 3: the Application (worked together, ~6 min) ────────────────────────
\boxguard[14]
\begin{notesbox}{Application: The Season Pass}
We will do this one together. A ski area records five things about each of its 3{,}400
season-pass holders: \textbf{pass number}, \textbf{age} (years), \textbf{pass tier}
(Gold / Silver / Bronze), \textbf{miles traveled} to the mountain, and \textbf{locker number}.

\begin{enumerate}[label=\textbf{\alph*.}, leftmargin=1.8em, itemsep=6pt]
  \item The individuals are \ans{the 3,400 season-pass holders}. Label each variable \textbf{C} or \textbf{Q}.

        \vspace{4pt}
        \renewcommand{\arraystretch}{1.3}
        \begin{tabularx}{\linewidth}{@{} Y Y Y Y Y @{}}
        \textbf{Pass \#} & \textbf{Age} & \textbf{Pass tier} & \textbf{Miles traveled} &
        \textbf{Locker \#} \\
        \ans{C} & \ans{Q} & \ans{C} & \ans{Q} & \ans{C} \\
        \end{tabularx}
  \item \textbf{What if we change the variable?} The resort stops recording age in years and
        records \textbf{age group} instead (Under 18 / 18--34 / 35--54 / 55+). Same people,
        same information --- which kind is it now, and what changed?
        \answerspace{1.8cm}{It is now categorical. Each holder carries a group label instead of
        a measured amount, so the resort can still count how many are in each band but can no
        longer average ages or say how far apart two holders are.}
  \item A report states: \textbf{``The average pass number is 5{,}120.''} Say in one sentence
        why that number tells you nothing about the pass holders.
        \answerspace{1.6cm}{Pass numbers are labels that identify holders, not amounts of
        anything, so averaging them produces a number that measures nothing.}
\end{enumerate}
\end{notesbox}

\end{document}
